\documentclass[11pt]{article}
\usepackage[T1]{fontenc}
\usepackage{lmodern,amsmath,amssymb,amsthm,mathtools}
\usepackage[margin=1in]{geometry}
\usepackage{microtype,booktabs,array,xcolor,enumitem}
\usepackage{fancyhdr,needspace}
\usepackage[colorlinks=true,linkcolor=blue!45!black,urlcolor=blue!45!black,citecolor=blue!45!black]{hyperref}
\usepackage{listings}
\hypersetup{pdftitle={AC-04: A certified upper-bound refinement},
pdfauthor={Research write-up prepared with ChatGPT},
pdfsubject={Partial result: asymptotic rank of the small Coppersmith-Winograd tensor},
pdfkeywords={AC-04, asymptotic rank, tensor, partial result, exact certificate}}
\setlength{\parindent}{0pt}
\setlength{\parskip}{6pt}
\setlength{\headheight}{15pt}
\pagestyle{fancy}\fancyhf{}
\fancyhead[L]{\small AC-04 \; / \; Asymptotic tensor rank}
\fancyhead[R]{\small Partial result, not a resolution}
\fancyfoot[C]{\thepage}
\newcommand{\AR}{\widetilde{R}}
\newcommand{\BR}{\underline{R}}
\newcommand{\C}{\mathbb C}
\newcommand{\Q}{\mathbb Q}
\newcommand{\rank}{\operatorname{rank}}
\newcommand{\diag}{\operatorname{diag}}
\newcommand{\Id}{\operatorname{id}}
\newcommand{\degto}{\trianglelefteq}
\newtheorem{theorem}{Theorem}[section]
\newtheorem{lemma}[theorem]{Lemma}
\newtheorem{proposition}[theorem]{Proposition}
\newtheorem{corollary}[theorem]{Corollary}
\theoremstyle{definition}\newtheorem{definition}[theorem]{Definition}
\theoremstyle{remark}\newtheorem{remark}[theorem]{Remark}
\lstset{basicstyle=\ttfamily\small,breaklines=true,columns=fullflexible,frame=single,rulecolor=\color{black!15}}
\title{\textbf{AC-04: A certified upper-bound refinement}\\[5pt]
\large Iterated slice extraction for the small Coppersmith--Winograd tensor}
\author{Research write-up prepared with ChatGPT}
\date{Request date: 13 September 2026}
\begin{document}
\begin{center}
{\LARGE\bfseries AC-04: A certified upper-bound refinement\\[7pt]}
{\large Iterated slice extraction for the small Coppersmith--Winograd tensor}\\[14pt]
{\small Research write-up prepared with ChatGPT \quad | \quad Request date: 13 September 2026}
\end{center}
\vspace{6pt}
\thispagestyle{fancy}
\begin{center}
\fcolorbox{black!30}{black!3}{\parbox{0.93\linewidth}{
\textbf{Scope and status.} This document does \emph{not} prove or disprove
$\AR(\mathrm{cw}_2)=3$. It derives the strict upper bound
$\AR(\mathrm{cw}_2)<3.923038$ from the speedup framework of Alman and Li,
with an analytic parameter argument and exact rational certificates.
The underlying speedup method is existing work. No claim of research priority,
independent peer review, or proof-assistant verification is made.}}
\end{center}

\begin{abstract}
For the tensor in repository problem AC-04, a rank-four normal form supplies a
full, isolated slice degeneration starting from the third tensor power. Repeated
squaring and re-extraction lead to a closed recurrence for dimensions and spectral
constraints. An analytic monotonicity lemma reduces the entire relevant spectral
parameter interval to the endpoint $2/3$. Three re-extractions, certified using
integer cube roots and rational interval arithmetic, give
$3\leq\AR(T)<3.923038$; the scalar endpoint is approximately
$3.923037967885872253$. The same recurrence, even with arbitrarily many iterations,
allows a formal scalar value of $3.9$. This last observation is an obstruction
\emph{only to deductions from these scalar constraints}, not a lower bound on
$\AR(T)$. The original equality question remains unproved in this work.
\end{abstract}

\section{The question and the result actually established}
The tensor is
\begin{equation}\label{eq:T}
T=\sum_{i=1}^{2}\bigl(e_0\otimes e_i\otimes e_i
+e_i\otimes e_0\otimes e_i+e_i\otimes e_i\otimes e_0\bigr)
\in\C^3\otimes\C^3\otimes\C^3.
\end{equation}
All tensor powers in this document are grouped, three-mode Kronecker powers:
the first factors are grouped together, and likewise the second and third.
The question is whether
\begin{equation}\label{eq:target}
\AR(T):=\lim_{k\to\infty}R(T^{\otimes k})^{1/k}=3.
\end{equation}
This is the exact target in AC-04 \cite{repo}, not the assertion that some fixed
power has rank or border rank $3^k$.

\begin{theorem}[Upper-bound refinement]\label{thm:main}
For the tensor \eqref{eq:T},
\[
\boxed{\qquad 3\leq\AR(T)<3.923038.\qquad}
\]
More precisely, $\AR(T)\leq\beta$, where $\beta$ is the unique nonnegative solution
of the scalar equation in Section~\ref{sec:certificate}, and
\[
3.92303796788587225299<\beta<3.92303796788587225301.
\]
\end{theorem}

The source benchmark is the bound $\AR(T)<3.931$ of Alman and Li
\cite{alman}. Their framework explicitly permits further iteration. The present
calculation changes the initial tensor power to three and makes the iteration
and endpoint certification explicit; it is not a new speedup theorem.
Their Table~1, using initial power four for the quoted second-method bound,
reports $3.930872$. This comparison is to the displayed benchmark, not an
assertion that all unpublished or subsequent optimizations have been ruled out.

The equality in \eqref{eq:target} is known to imply $\omega=2$
\cite[Theorem 1.2]{alman}. We do not use that implication to dismiss the problem,
and we do not assert its converse. The stronger equality is simply not a
conclusion of the proof below.

\section{An exact rank-four normal form}
Let
\[
P=\sum_{\sigma\in S_3}
 a_{\sigma(1)}\otimes a_{\sigma(2)}\otimes a_{\sigma(3)},
\]
where $a_1,a_2,a_3$ is the standard basis of $\C^3$.
Set $w_1=e_0$, $w_2=e_1+\mathrm i e_2$, and $w_3=e_1-\mathrm i e_2$.
Then
\[
w_2\otimes w_3+w_3\otimes w_2
=2(e_1\otimes e_1+e_2\otimes e_2),
\]
so, for the invertible map $L a_i=w_i$,
\begin{equation}\label{eq:normal}
T=\tfrac12L^{\otimes3}P.
\end{equation}
Nonzero scalar factors and invertible changes of basis preserve rank, border
rank, and asymptotic rank. It is therefore enough to work with $P$.

Define
\[
v_0=(1,1,1)^\mathsf T,\quad v_1=(1,-1,-1)^\mathsf T,\quad
v_2=(-1,1,-1)^\mathsf T,\quad v_3=(-1,-1,1)^\mathsf T.
\]
A coefficient check gives the exact decomposition
\begin{equation}\label{eq:rankfour}
P=\frac14\sum_{\ell=0}^3v_\ell\otimes v_\ell\otimes v_\ell.
\end{equation}
Indeed, the coefficient is one when all three indices differ, and zero
otherwise. In particular, $R(T)\leq4$.

The three slices of a flattening of $P$ are linearly independent: each has its
nonzero entries on a different pair of off-diagonal positions. Thus every
flattening of $P$ has rank three. Flattening commutes with the grouped tensor
product, and matrix ranks multiply. Consequently,
\[
R(T^{\otimes k})\geq3^k\quad(k\geq1).
\]
Submultiplicativity of tensor rank makes $\log R(T^{\otimes k})$ subadditive;
its normalized limit is its infimum. This proves both existence of the
asymptotic-rank limit and the lower bound in Theorem~\ref{thm:main}.

For extraction, let $V=[v_0\;v_1\;v_2\;v_3]$ and
$D=\diag(1,1,-1,-1)$. Direct multiplication gives
\begin{equation}\label{eq:sign}
VDV^\mathsf T=
\begin{pmatrix}0&0&0\\0&0&4\\0&4&0\end{pmatrix}=XY^\mathsf T,
\qquad
X=\begin{pmatrix}0&0\\4&0\\0&4\end{pmatrix},\quad
Y=\begin{pmatrix}0&0\\0&1\\1&0\end{pmatrix}.
\end{equation}
All four diagonal coefficients in $D$ are nonzero, while the contracted image
has rank two. This is the finite algebra needed below.

\section{The extraction mechanism and its hypotheses}
We spell out the relevant mechanism so the numerical recurrence is not treated
as a free-standing conjecture. It is the one-slice specialization of the
free-lunch construction in \cite[Section 5]{alman}.

Write
\[
D_r=\sum_{i=1}^r a_i\otimes b_i\otimes c_i,
\qquad L_s=\sum_{i=1}^s u_i\otimes v_i\otimes w.
\]
Thus $D_r$ is diagonal and $L_s$ has dimensions $(s,s,1)$; in matrix multiplication
notation it is $\langle1,s,1\rangle$. Put $L_0=0$.
The notation $K\degto S$ means that $K$ is a degeneration of $S$, with maps
allowed to have rational-function coefficients in a degeneration parameter.

\subsection{A kernel-compression lemma}
\begin{lemma}\label{lem:extract}
Suppose maps $A,B,C$ send a tensor $S$ to $K$, and the first two target dimensions
are $d_A,d_B$. Let $f$ be a functional on the third source space such that
\[
(A\otimes B\otimes f)S=0.
\]
If the source contraction $H=(\Id\otimes\Id\otimes f)S$ has matrix rank $M$,
then $S$ degenerates to $K\oplus L_t$ for every positive integer
$t\leq M-d_A-d_B$.
The new slice is isolated in its third mode. Its source fullness index is $M$.
The statement also applies over a rational-function field and can be composed
with a degeneration to $K$.
\end{lemma}
\begin{proof}
Regard rows of $A,B$ as functionals on the respective source spaces. Let
\[
E=\{u:uHB^\mathsf T=0\},\qquad
F=\{v:AHv^\mathsf T=0\}.
\]
Their codimensions are at most $d_B,d_A$. Restricting a matrix to a subspace of
codimension $c$ on either side decreases its rank by at most $c$. Hence the
restriction of $H$ to $E\times F$ has rank at least $M-d_A-d_B$.
Choose row maps $A',B'$ from these subspaces and then restrict and normalize a
rank-$t$ submatrix so that
$(A'\otimes B'\otimes f)S=L_t$.

Augment the three maps to $(A,A')$, $(B,B')$, and $(C,f)$. There are eight
blocks indexed by $1,2$ in the three modes. The $111$ block is $K$ and the
$222$ block is $L_t$. By construction, blocks $112,122,212$ vanish identically.
Assign weights $(1,\lambda)$ to the two first-mode blocks, the same weights
to the second-mode blocks, and weights $(\lambda^2,1)$ to the third-mode
blocks. Divide the resulting tensor by $\lambda^2$.
The $111$ and $222$ blocks survive. Blocks $121$ and $211$ have an extra factor
$\lambda$, and block $221$ an extra factor $\lambda^2$; they disappear.
This proves the degeneration.

Every block ending in the new third-mode coordinate, other than $222$, is
identically zero even before taking the limit. This is isolation. The
functional selecting that coordinate is a nonzero scalar multiple of $f$, so
its contraction on the \emph{source} still has rank $M$; this is fullness index
$M$.

For rational-function input maps, perform these linear-algebra operations over
that field. Normalize the extracted rank-$t$ matrix before taking the limit,
so that rank is not lost in its leading coefficient. The finitely many entries
of all maps have finite pole orders. Replacing the new parameter by a
sufficiently high power of the original one makes every unwanted block tend
to zero. The desired $K$ block retains its original limit, and the normalized
slice stays constant. Isolation remains exact, and source contraction rank is
unchanged over the rational-function field.
\end{proof}

\subsection{Why full and isolated are essential}
For a degeneration $K\oplus L_t\degto S$, with maps $A_\lambda,B_\lambda,C_\lambda$,
call the distinguished $L_t$ \emph{isolated} if the third-mode coordinate
selecting it has identically zero coefficients outside its own first/second-mode
block before the limit. Call it \emph{full} when the associated contraction of
$S$ has rank equal to each of the first two source dimensions.
These are properties of the maps defining the degeneration, not merely of the
limiting direct sum.

Tensor products preserve both properties: the relevant source contraction is a
Kronecker product, so ranks multiply; and an off-block contribution to the
product's distinguished third coordinate contains an identically zero factor.
This is the mechanism also formalized in \cite[Propositions 5.1--5.6]{alman}.
Deleting a full, isolated slice leaves its source functional satisfying the
zero-contraction hypothesis of Lemma~\ref{lem:extract}. Without isolation, an
$O(\lambda)$ cross-block term would not suffice for this argument.

\section{The base degeneration: dimensions 27, 18, and 72}\label{sec:base}
Use the rank-four decomposition \eqref{eq:rankfour} three times. Thus
$P^{\otimes3}$ is a restriction of $D_{64}$ with first two maps
$A=V^{\otimes3}$ and third map $4^{-3}V^{\otimes3}$.
The source third-mode functional with coefficient matrix $D^{\otimes3}$ has
64 nonzero entries, and its image contraction is
\[
A D^{\otimes3}A^\mathsf T
=(VDV^\mathsf T)^{\otimes3}
=X^{\otimes3}(Y^{\otimes3})^\mathsf T,
\]
of rank $2^3=8$.

Adjoin $L_8$ to the source. Map its first two modes by $X^{\otimes3}$ and
$Y^{\otimes3}$, and its third mode by zero for the map onto $P^{\otimes3}$.
Extend the chosen third functional with coefficient $-1$ on the $L_8$ slice.
The contracted image cancels exactly. On the source, the contraction matrix is
\[
H=D^{\otimes3}\oplus(-I_8),
\]
which is nonsingular of size 72. Lemma~\ref{lem:extract}, with $d_A=d_B=27$,
yields a full and isolated degeneration
\begin{equation}\label{eq:base}
U_0:=P^{\otimes3}\oplus L_{18}\degto S:=D_{64}\oplus L_8,
\qquad 18=72-2\cdot27.
\end{equation}
It holds in all three slice orientations, because $P$ is symmetric in its modes.

The file \texttt{code/algebra\_check.py} checks the complete finite linear
algebra here using exact SymPy matrices. The two annihilator spaces have
dimension 45. Their contracted $45\times45$ matrix has rank 18; the certificate
supplies an $18\times18$ minor with determinant $-1$. Thus the guaranteed slice
size is realized by this explicit base construction, not inferred from a
floating-point numerical rank.

\Needspace{18\baselineskip}
\section{Repeated squaring and re-extraction}\label{sec:iteration}
\begin{proposition}\label{prop:recurrence}
There are degenerations
\[
U_j=K_j\oplus L_{t_j}\degto S^{\otimes2^j}\qquad(j\geq0),
\]
whose distinguished slices are full and isolated. The first two dimensions of
$K_j$ are $d_j$, and those of the source are $M_j$, where
\begin{align}\label{eq:dimensions}
(d_0,t_0,M_0)&=(27,18,72),\\
d_{j+1}&=d_j^2+2d_jt_j,\nonumber\\
t_{j+1}&=t_j^2+2d_j^2,\nonumber\\
M_{j+1}&=M_j^2=t_{j+1}+2d_{j+1}.\nonumber
\end{align}
The third dimension of $K_j$ need not equal $d_j$.
\end{proposition}
\begin{proof}
The base case is \eqref{eq:base}. Square the degeneration at step $j$ and
separate its distinguished summand $L_{t_j^2}$. The remaining tensor is
\[
K_{j+1}=K_j^{\otimes2}\oplus(K_j\otimes L_{t_j})
\oplus(L_{t_j}\otimes K_j),
\]
whose first two dimensions are $d_j^2+2d_jt_j$.
The removed slice has source fullness $M_j^2$ and is isolated. Its source
functional therefore meets the cancellation condition after its image
variables are removed. Re-extract with Lemma~\ref{lem:extract} to obtain
\begin{align*}
t_{j+1}&=M_j^2-2(d_j^2+2d_jt_j)\\
&=(t_j+2d_j)^2-2d_j^2-4d_jt_j=t_j^2+2d_j^2.
\end{align*}
The new slice is again full and isolated. This proves the induction, including
$M_{j+1}=t_{j+1}+2d_{j+1}$.
\end{proof}

This is the repeated form of the mechanism in \cite[Theorem 6.2]{alman}.
There is a useful exact closed form:
\begin{equation}\label{eq:closed}
d_j=\frac{72^{2^j}-(-9)^{2^j}}3,\qquad
t_j=\frac{72^{2^j}+2(-9)^{2^j}}3,\qquad M_j=72^{2^j}.
\end{equation}
Indeed, both $2d_j+t_j$ and $t_j-d_j$ are squared by the recurrence, with initial
values 72 and $-9$. The first values are
\begin{center}\small
\begin{tabular}{rrrr}
\toprule
$j$&$d_j$&$t_j$&$M_j$\\\midrule
0&27&18&72\\
1&1,701&1,782&5,184\\
2&8,955,765&8,962,326&26,873,856\\
3&240,734,697,754,005&240,734,740,800,726&722,204,136,308,736\\
\bottomrule
\end{tabular}
\end{center}
The large dimensions are tracked symbolically; no tensor of these dimensions
is explicitly allocated in the computation.

\Needspace{10\baselineskip}
\section{Turning the degenerations into scalar bounds}\label{sec:spectral}
We use Strassen's asymptotic-spectrum duality in its tensor form, as stated and
applied in \cite[Section 4]{alman}. Precisely, for complex tensors there is a
family of spectral points $\phi$ such that
\[
\AR(P)=\max_\phi\phi(P).
\]
A spectral point is additive for direct sums, multiplicative for grouped tensor
products, monotone under degeneration, and normalized by $\phi(D_r)=r$.
For the three orientations of $L_s$, its values are $s^{\theta_1}$,
$s^{\theta_2}$, $s^{\theta_3}$, with
\begin{equation}\label{eq:theta}
0\leq\theta_i\leq1,\qquad\theta_1+\theta_2+\theta_3\geq2.
\end{equation}
These are standard external theorems used here, not assertions that the
attached code proves. In particular, at least one orientation has
$\theta\in[2/3,1]$.

Fix a spectral point and choose that orientation. Write
\[
z=\phi(P),\qquad x=z^3,\qquad h(\theta)=64+8^\theta.
\]
The exact rank-four upper bound gives $0\leq z\leq4$, hence $0\leq x\leq64$.
Define
\begin{align}\label{eq:p}
p_0(x,\theta)&=x+18^\theta,\\
p_{j+1}(x,\theta)&=p_j(x,\theta)^2-t_j^{2\theta}+t_{j+1}^{\theta}.
\nonumber
\end{align}
The subtraction removes precisely the square of the old distinguished slice,
which is replaced by the new one. Additivity and multiplicativity give
$\phi(U_j)=p_j(x,\theta)$. Consequently every spectral point obeys
\begin{equation}\label{eq:spectralconstraint}
p_j(x,\theta)\leq h(\theta)^{2^j}\qquad(j\geq0).
\end{equation}
Although the displayed recursion subtracts a constant, $p_j$ is the value of an
actual direct sum: as a polynomial in $x$, all its coefficients are
nonnegative, its constant term is $t_j^\theta$, and it is strictly increasing
for $x\geq0$.

\subsection{An analytic endpoint lemma}
Numerically searching for the worst $\theta$ would not certify an interval of
parameters. The next lemma eliminates that search.

\begin{lemma}\label{lem:endpoint}
For fixed $x\in[0,64]$ and each $j\geq0$, the function
\[
q_j(x,\theta)=\frac{p_j(x,\theta)}{h(\theta)^{2^j}}
\]
is strictly increasing for $\theta>0$.
\end{lemma}
\begin{proof}
For $j=0$, the numerator of $\partial q_0/\partial\theta$, after multiplying
by the positive denominator $h^2$, is
\[
64\,18^\theta\log18-x\,8^\theta\log8
+18^\theta8^\theta\log(18/8).
\]
It is strictly positive: $x\leq64$, $18>8$, and $\theta>0$.

Normalize the recursion to obtain
\begin{equation}\label{eq:q}
q_{j+1}=q_j^2+E_j(\theta),\qquad
E_j(\theta)=\frac{t_{j+1}^{\theta}-t_j^{2\theta}}
                 {h(\theta)^{2^{j+1}}}.
\end{equation}
We have $t_{j+1}>t_j^2$ and, inductively from $t_0=18>8$,
$t_j>8^{2^j}$. If $A>B>0$ and $\theta>0$, then
\[
\frac{d}{d\theta}\log(A^\theta-B^\theta)
=\log B+\frac{A^\theta\log(A/B)}{A^\theta-B^\theta}>\log B.
\]
On the other hand, $h'/h<\log8$. Apply these inequalities with
$A=t_{j+1}$ and $B=t_j^2$. They imply
\[
\frac{E'_j(\theta)}{E_j(\theta)}
>2\log t_j-2^{j+1}\log8>0.
\]
Thus $E_j$ is positive and strictly increasing. Since $q_j$ is positive,
\eqref{eq:q} proves the claim by induction.
\end{proof}

For the chosen orientation $\theta\geq2/3$, equations
\eqref{eq:spectralconstraint} and Lemma~\ref{lem:endpoint} imply
\begin{equation}\label{eq:endpointconstraint}
p_j(z^3,2/3)\leq68^{2^j}.
\end{equation}
This is an \emph{upper} bound: monotonicity in $\theta$ makes the endpoint
constraint weaker, so every actual spectral point must satisfy it.
No permutation-invariance of $\phi$ was assumed; the symmetry of $P$ supplies
three separate oriented degenerations.

\section{Exact certificate for the strict bound}\label{sec:certificate}
Define $\beta$ by
\begin{equation}\label{eq:beta}
p_3(\beta^3,2/3)=68^8.
\end{equation}
The left side is strictly increasing in $\beta\geq0$. The certified bracket
below proves existence and uniqueness of the nonnegative root.
Every spectral value $z$ satisfies \eqref{eq:endpointconstraint}, so
$z\leq\beta$. Duality gives $\AR(P)\leq\beta$.

Only cube roots of positive integers are needed for evaluation, because
$t^{2/3}=\sqrt[3]{t^2}$. For any nonnegative integer $N$, let
\[
a=\left\lfloor\sqrt[3]{N\,10^{120}}\right\rfloor.
\]
Integer comparisons determine $a$ exactly. Then
\[
\frac{a}{10^{40}}\leq\sqrt[3]{N}\leq\frac{a+1}{10^{40}},
\]
with equality handled exactly for perfect cubes. Interval addition,
subtraction, squaring, and division by positive rational numbers preserve
rigorous enclosures.

The standard-library checker \texttt{code/certify.py} propagates these intervals
through \eqref{eq:p}. At the rational number $\alpha=3.923038$, it proves
\[
\frac{p_3(\alpha^3,2/3)}{68^8}>1.
\]
For a readable indication of the margin, the enclosed value is approximately
\[
1.0000001724140039588978237928635588.
\]
The comparison with one is made using exact rational endpoints, not this
rounded decimal. It follows that $\beta<3.923038$.

The same checker independently proves the two strict inequalities
\begin{align*}
p_3\bigl((3.92303796788587225299)^3,2/3\bigr)&<68^8,\\
p_3\bigl((3.92303796788587225301)^3,2/3\bigr)&>68^8.
\end{align*}
Both inputs here are terminating rationals. These checks give the bracket in
Theorem~\ref{thm:main}. Combining it with the flattening lower bound completes
the proof of that theorem.

\subsection{Numerical progression and its proper interpretation}
For context, the scalar roots for successive values of $j$ are shown below.
Rows not used in the exact certificate are exploratory high-precision
evaluations, not additional computer-assisted theorems.
\begin{center}
\begin{tabular}{rl}
\toprule
Re-extractions $j$ & Initial power $n=3$: endpoint root\\\midrule
0 & $3.939328446458942411\ldots$\\
1 & $3.923413930957849239\ldots$\\
2 & $3.923038786084391559\ldots$\\
3 & $3.923037967885872253\ldots$ \quad(exact bracket certified)\\
4 & $3.923037967878664608\ldots$\\
5 & $3.923037967878664608\ldots$\\\bottomrule
\end{tabular}
\end{center}
For comparison, starting with $n=4$ gives endpoint values approximately
$3.9334837730$ at $j=0$ and $3.9308711169$ at $j=1$, consistent with the
upward-rounded benchmark in \cite[Table 1]{alman}. The initial power that is
best before re-extraction need not remain best after it. No global optimality
over initial powers, mixing patterns, or other degeneration methods is claimed.

The step $j=3$ involves a $24$th power of $P$ within a much larger direct sum.
This does \emph{not} provide an explicit assertion
$R(P^{\otimes24})\leq3.923038^{24}$. The argument passes from direct-sum
constraints to asymptotic rank through spectral duality. Confusing these two
claims would invalidate the interpretation of the result.

\section{A rigorous limitation of this particular recurrence}\label{sec:barrier}
The rapid apparent numerical stabilization has an analytic counterpart.
We give a deliberately modest threshold, $3.9$, for which a simple uniform
argument works.

\begin{proposition}[Method-local scalar obstruction]\label{prop:barrier}
The constraints \eqref{eq:spectralconstraint} for all $j\geq0$, generated by the
$n=3$ binary re-extraction family in all three orientations and supplemented by
\eqref{eq:theta} and $0\leq z\leq4$, admit the formal assignment
\[
z=39/10,\qquad \theta_1=\theta_2=\theta_3=2/3.
\]
Therefore these constraints alone cannot imply $z<3.9$, and in particular
cannot establish the target $\AR(P)=3$.
\end{proposition}
\begin{proof}
Put $x=(39/10)^3$, $h=68$, and $c=39/40$. Write $q_j=q_j(x,2/3)$.
The exact checker verifies
\[
q_0=\frac{(39/10)^3+18^{2/3}}{68}<c,
\qquad
E_0=\frac{1782^{2/3}-18^{4/3}}{68^2}<c-c^2=\frac{39}{1600}.
\]
For reference, these values are approximately $0.9733424332$ and
$0.02158547644$. The exact comparisons, not the decimal displays, are used.

For $j\geq1$, formula \eqref{eq:closed} gives $t_{j+1}\leq72^{2^{j+1}}$, hence
\begin{align*}
0<E_j(2/3)
&<\frac{t_{j+1}^{2/3}}{68^{2^{j+1}}}
\leq\left(\frac{72^{2/3}}{68}\right)^{2^{j+1}}\\
&<\left(\frac9{34}\right)^4<\frac{39}{1600}.
\end{align*}
Here $72^{2/3}<18$ follows from the integer inequality $72^2<18^3$.
The last comparison is rational and exact. If $q_j<c$, the recurrence gives
$q_{j+1}=q_j^2+E_j<c^2+(c-c^2)=c$. Thus $q_j<c<1$ for every $j$.
All orientations have the same chosen exponent, so all displayed constraints
are satisfied, as are $\theta_1+\theta_2+\theta_3=2$ and $z\leq4$.
\end{proof}

\textbf{This is not a tensor lower bound.} We have not constructed a spectral
point with $\phi(P)=3.9$. A feasible assignment for a subset of the necessary
spectral inequalities need not extend to a genuine spectral point on all
tensors. Accordingly, Proposition~\ref{prop:barrier} says nothing of the form
$\AR(P)\geq3.9$. Its use is to identify a precise missing ingredient: more
iterations of these same scalar inequalities cannot close the gap to three.

\section{What remains unproved, and two structural reformulations}
A complete positive solution still needs to establish the quantified statement
\begin{equation}\label{eq:quantifier}
\text{for every }\varepsilon>0\text{ there exists }k\geq1\text{ such that }
R(T^{\otimes k})<(3+\varepsilon)^k.
\end{equation}
By submultiplicativity this is equivalent to \eqref{eq:target}. Equivalently,
one needs the upper growth estimate $R(T^{\otimes k})=3^{k+o(k)}$.
Replacing rank by border rank leads to the same asymptotic target, using the
standard rank/border-rank asymptotic relation discussed in \cite{alman}.
No such estimate is supplied here.

Two elementary descriptions locate the issue without asserting that it is
solved. First, identify the three indices with the three nonzero elements of
$G=(\mathbb Z/2\mathbb Z)^2$. Then the support of $P$ is exactly
\[
\{(a,b,c)\in(G\setminus\{0\})^3:a+b+c=0\}.
\]
Thus $P^{\otimes k}$ is a restricted group-convolution tensor. The ambient group
has $4^k$ elements while each restricted mode has $3^k$ elements. An algorithm
with $3^{k+o(k)}$ bilinear multiplications for this restricted convolution would
prove the desired upper bound. The ambient Fourier decomposition merely
recovers the four-term construction; it does not provide that improvement.
This support identification follows by checking the three possible nonzero
values in one coordinate.

Second, index a cyclic group by $0,1,2$ and write
\[
C_3=\sum_{a+b+c=0\;({\rm mod}\,3)}e_a\otimes e_b\otimes e_c.
\]
Its nine terms consist of the six terms of $P$ (with indices relabeled) and
three diagonal terms. Hence
\[
P=C_3-D_3.
\]
A Fourier decomposition gives $R(C_3)=3$. It is not valid to infer
$\AR(C_3-D_3)\leq3$ from this identity: subtraction is not a tensor restriction,
and asymptotic rank has no subtraction rule.

There is even a small obstruction to a tempting diagonal degeneration of
this identity. Suppose local monomial weights $a_i,b_j,c_k$ preserve all six
permutation terms at a common total weight $m$. Swapping indices in those
terms forces $a_i-b_i$ to be constant, and likewise $a_i-c_i$ to be constant.
The weights of the three unwanted diagonal terms then have sum $3m$.
They cannot all be strictly greater than $m$, so this local weighting alone
cannot kill all three diagonal terms while preserving all six desired ones.
This observation excludes only this elementary diagonal-weighting attempt;
it does not exclude general degenerations, restrictions of larger powers,
or a positive solution to AC-04.

Finally, the known finite result $\BR(T^{\otimes2})=16$ \cite{conner} cannot
be promoted to $\AR(T)=4$. The asymptotic rank is an infimum over all powers.
Neither this finite lower bound nor the bound proved in this report resolves
which point in the surviving interval is the asymptotic rank.

\section{Audit, reproducibility, and delivery scope}
The theorem proof uses an existing asymptotic-spectrum theorem and the
slice-extraction mechanism, with their exact roles stated in
Sections~3--6. The new calculations are the initial parameter choice,
explicit repeated recurrence, all-parameter endpoint proof, exact bound
certificate, and the limited scalar obstruction. Priority for these refinements
has not been established by the literature search.

The scalar certificate uses only the Python standard library. The separate
finite-matrix audit requires SymPy; it was run with version 1.14.0.
The following commands reproduce the supplied checks from the package root:
\begin{lstlisting}
python3 code/certify.py --output results/exact_certificate.json
python3 code/algebra_check.py --output results/algebra_certificate.json
(cd code && python3 -m unittest -v test_certify.py)
\end{lstlisting}
The delivered run passes eight regression tests and both exact certificate
programs. The main numeric test uses no floating-point comparison. The
high-precision exploratory table is generated separately and is not used as
a proof of a strict inequality.

No proof assistant was run. No independent referee has audited the full
argument. Passing the algebra and interval checks is not equivalent to either
formal verification or outside peer review. The local status of this package
is \textbf{partial result}; the repository itself was not edited, and this
package must not be reported as a solution or disproof of AC-04. The repository's
resolution guidance likewise distinguishes supporting results from a
complete resolution of the original target \cite{contributing}.

\clearpage
\begin{thebibliography}{9}
\bibitem{repo}
OpenProblemsInNLA, \emph{AC-04: Asymptotic rank of the small
Coppersmith--Winograd tensor}, canonical problem entry, consulted for this
request.\newline
\url{https://github.com/ajt60gaibb/OpenProblemsInNLA/tree/main/arithmetic-and-complexity/AC-04}

\bibitem{alman}
J. Alman and B. Li, \emph{Asymptotic Rank Speedup Theorems, Revisited}.
41st Computational Complexity Conference (CCC 2026), LIPIcs 383,
36:1--36:42, published 23 July 2026.
DOI: \href{https://doi.org/10.4230/LIPIcs.CCC.2026.36}{10.4230/LIPIcs.CCC.2026.36}.
Detailed theorem locators in this report refer to the arXiv v1 exposition:
Theorem 4.1 and Propositions 4.1--4.2 (spectrum),
Sections 5.2--5.3 (extraction, fullness and isolation),
Theorems 6.1--6.2 (speedups), and Section 7.1/Table 1 (CW bounds).\newline
\url{https://arxiv.org/html/2605.21738v1}

\bibitem{conner}
A. Conner, H. Huang, and J. M. Landsberg,
\emph{Bad and good news for Strassen's laser method: Border rank of the
$3\times3$ permanent and strict submultiplicativity},
arXiv:2009.11391 (2020), Theorem 1.1.\newline
\url{https://arxiv.org/abs/2009.11391}

\bibitem{contributing}
OpenProblemsInNLA, \emph{Contributing}, especially ``Reporting a resolution''
and the distinction between supporting lemmas and complete target verification.\newline
\url{https://github.com/ajt60gaibb/OpenProblemsInNLA/blob/main/CONTRIBUTING.md}
\end{thebibliography}
\end{document}
